\section{Related Work}
\paragraph{VLM bias measurement and post-hoc editing.}
Audits have documented demographic and stereotype-related disparities in VLM classification and retrieval \citep{agarwal2021evaluating,hall2023gender,zhou2022vlstereoset,hamidieh2024implicit}. Prompt and projection methods remove directions specified by biased prompt pairs \citep{berg2022prompt,chuang2023debiasing}; BEND-VLM uses attribute-labeled references for query-specific correction \citep{gerych2024bendvlm}; SANER learns a societal-attribute neutralizer from a predefined lexicon without per-image attribute labels \citep{hirota2025saner}; and recent work treats bias as an attribute-predictive subspace \citep{jung2024unified,zhao2026subspace}. These methods start from an attribute or concept. We instead start from within-class failure and test post hoc whether the resulting editor improves group metrics.

\paragraph{Robustness without training-group labels.}
Group DRO directly optimizes worst-group risk when group annotations are available \citep{sagawa2020groupdro}. JTT uses early errors to upweight hard examples \citep{liu2021jtt}, EIIL infers environments from loss \citep{creager2021eiil}, and last-layer methods improve robustness with limited group supervision \citep{kirichenko2023dfr,qiu2023afr}. Correct-N-Contrast forms contrastive sets from error-derived groups \citep{zhang2022correct}. Several methods avoid group labels during fitting but still use group-labeled validation data for hyperparameter or checkpoint selection; we distinguish this regime from both fully group-free selection and attribute-supervised adaptation in Table~\ref{tab:classification}. Our setting is more restrictive in model access: image and text encoders remain frozen, and intervention occurs in their shared embedding space. It is also diagnostically different: we do not assume that errors recover the true groups, and our CelebA result demonstrates the cost of that assumption.

\paragraph{Representation subspaces.}
INLP iteratively removes linearly decodable information \citep{ravfogel2020inlp}, and LEACE gives a closed-form linear concept-erasure solution \citep{belrose2023leace}. Recent VLM work also argues that social bias is distributed across subspaces rather than individual coordinates \citep{zhao2026subspace}. We evaluate both a transparent mean-contrast direction and iterative risk probes. On Waterbirds, mean contrast outperforms the rank-4 construction; on CelebA, even a highly risk-separable rank-1 probe harms WGA. Our focus is therefore not whether risk is decodable, but whether suppressing it preserves task semantics and improves held-out groups.

\paragraph{Difference from supervised concept erasure.}
Concept erasure begins with a named concept and learns how to remove it from a representation. INLP and LEACE therefore answer: how can a known concept be made linearly inaccessible? \textsc{LatentRisk-VLM} begins instead with model failures and asks which latent directions are repeatedly activated by those failures. The resulting direction may overlap with a spurious attribute, but no attribute identity is assumed or required. We do not claim to recover the true sensitive attribute; a risk direction is only a failure-related feature pattern. If it aligns with the annotated attribute, failure supervision may approximate attribute-supervised erasure. If it only partially aligns, it may remove the failure-relevant component while retaining other task semantics. If alignment is low, it may capture factors beyond the named spurious attribute.
